\section{A necessary quadratic limit}\label{sec:limit}

The power maps allow the fixed twist by $P$ to disappear in a limit.
We prove that the quadratic structure survives at every base point,
even where a raw pointwise limit need not be plurisubharmonic.
The argument has four steps: trace the dual norms, normalize and
extract their coefficient limit, prove that the limit is nonzero,
and recover its intrinsic plurisubharmonic representative.
We establish this statement for an arbitrary vector bundle on $S$.

\begin{proposition}\label{prop:limit}
Let $G$ be a holomorphic vector bundle of positive rank on
$S=\PP^1\times\PP^1$, and put $F=G^\vee$.  Suppose that
$f_m^*G\otimes P$ admits a smooth strictly Griffiths-positive Hermitian
metric for infinitely many positive integers $m$.  Then there is a
function
\[
q:F\longrightarrow [0,\infty)
\]
which is plurisubharmonic on the total space, is not identically zero,
and has the following property: for every $x\in S$ there is a finite
positive semidefinite Hermitian form $H_x$ on $F_x$ such that
\[
q(x,v)=H_x(v,v)\qquad(v\in F_x).
\]
In particular, this fiberwise assertion holds at every base point.
\end{proposition}

\begin{proof}
Put $r=\operatorname{rank}G$.  Choose an unbounded sequence of the
indicated integers and let
$\|\cdot\|_m$ denote the smooth dual metric on
\[
(f_m^*G\otimes P)^\vee=f_m^*F\otimes P^{-1}.
\]
Its squared norm is plurisubharmonic on the total space.  We recall the
sign in this standard implication.  For a smooth Hermitian metric with
matrix $h$, choose a holomorphic frame normal at the point in question,
so that $h=\Id$ and its first derivatives vanish there.  Along a
holomorphic disk with coordinate $\zeta$, write the Hermitian curvature
as $iC\,d\zeta\wedge d\bar\zeta$.  In the convention used here,
\[
C=-\frac{\partial^2h}{\partial\zeta\,\partial\bar\zeta}
\]
at the normal point.  Thus, for a holomorphic section $s$ along the disk,
\begin{equation}\label{eq:limit-dual-hessian}
\frac{\partial^2}{\partial\zeta\,\partial\bar\zeta}\|s\|^2
=\|D'_{\partial/\partial\zeta}s\|^2-\langle Cs,s\rangle.
\end{equation}
Duality reverses Griffiths positivity, so the second term on the right
is nonnegative for the dual metric.  Applying this calculation to the
section of the pulled-back bundle defined by any disk in its total
space proves the assertion. For this classical norm criterion and
duality, see \cite[Remark~1.3 and Lemma~1.7]{Umemura1973} and
\cite[Section~2]{Raufi2015}; the same implication is recorded in
\cite[Proposition~2.6]{ForstnericKusakabe}.

\smallskip\noindent\textbf{Trace quadratics and normalization.}
Fix, for each $m$, the holomorphic isomorphism
$f_m^*P^{-1}\simeq P^{-m}$.  If $t$ is a nonvanishing holomorphic frame
of $P^{-1}$ on an open set $U\subset S$, define on $F|_U$
\begin{equation}\label{eq:limit-trace}
q_{m,t}(x,v)
=\frac{1}{m^2}\sum_{y\in f_m^{-1}(x)}
\bigl\|v\otimes(f_m^*t)^{1/m}(y)\bigr\|_{m,y}^2,
\end{equation}
where the sum counts multiplicities.  The root in this expression is
local on the source.  Indeed, in a local frame $s$ of $P^{-1}$ the
section $f_m^*t$ has the form $\eta(y)s^m$, with $\eta$ holomorphic and
nowhere zero.  A local holomorphic logarithm of $\eta$ gives the root
$\exp((\log\eta)/m)s$.  Two choices differ by an $m$-th root of unity,
which does not change the squared norm.  The expression
\eqref{eq:limit-trace} is consequently well defined and is a positive
definite Hermitian quadratic in $v$.

We verify its regularity through the branch locus.  Off that locus,
every local inverse branch of $f_m$ gives a holomorphic map
\[
(x,v)\longmapsto
\bigl(y(x),v\otimes(f_m^*t)^{1/m}(y(x))\bigr)
\]
into $f_m^*F\otimes P^{-1}$.  Hence each summand, and their sum, is
plurisubharmonic.  The unordered multiset of roots of $y^m=x$ varies
continuously with $x$, counting coalescing roots with their
multiplicities.  Applied in the two factors, and using local source
frames for the norm in \eqref{eq:limit-trace}, this proves continuity
of the sum across the branch locus, including a crossing of its two
components.

The continuous trace is plurisubharmonic across the branch locus as
well. In a local total-space chart, the inverse image of the branch locus
under the bundle projection $F|_U\to U$ is an analytic set, contained in
$\{z_1z_2=0\}$ after choosing affine
coordinates on the two base factors. It is therefore pluripolar.
The trace is locally bounded above there by continuity, so the
removable-singularities theorem for plurisubharmonic functions
\cite[Chapter~I, Theorem~5.24]{DemaillyCADG} gives a unique
plurisubharmonic extension. Its upper-limit values agree with the
continuous trace. This proves joint plurisubharmonicity of $q_{m,t}$
on all of $F|_U$.

If $t'=gt$ is another nonvanishing frame, then
\begin{equation}\label{eq:limit-frame-law}
q_{m,t'}(x,v)=|g(x)|^{2/m}q_{m,t}(x,v).
\end{equation}
Choose smooth background Hermitian norms $|\cdot|_0$ on $F$ and
$P^{-1}$.  By \eqref{eq:limit-frame-law}, the continuous quadratic
\[
N_m(x,v)=q_{m,t}(x,v)|t(x)|_0^{-2/m}
\]
is globally defined. We use $N_m$ to normalize the traces globally;
the local functions $q_{m,t}$ retain the plurisubharmonicity used
below. The background unit sphere bundle is compact,
and $N_m$ is positive there.  Multiplying the dual metric's squared
norm by a positive constant, we may therefore arrange
\begin{equation}\label{eq:limit-normalization}
\max_{x\in S,\ |v|_0=1}N_m(x,v)=1.
\end{equation}
This rescaling does not change the curvature.  Homogeneity gives
\begin{equation}\label{eq:limit-local-bound}
0\le q_{m,t}(x,v)\le |t(x)|_0^{2/m}|v|_0^2.
\end{equation}

\smallskip\noindent\textbf{A measurable coefficient limit.}
In a holomorphic trivialization of $F|_U$, write
\[
q_{m,t}(x,v)=v^*H_m(x)v.
\]
On each relatively compact subset of $U$, the positive matrices $H_m$
are uniformly bounded, by \eqref{eq:limit-local-bound}.  Weak compactness
in $L^2$ and a diagonal subsequence on a countable collection of charts
and compact exhaustions give distributional limits $H$ for all the
coefficient matrices.  These limits belong to $L^\infty_{\mathrm{loc}}$,
are Hermitian, and are positive semidefinite almost everywhere.  For
example, positivity and the same local upper bound pass to
$c^*Hc$ against every nonnegative test function for each
$c\in(\mathbb Q+i\mathbb Q)^r$; density then gives the matrix
inequalities on a single set of full measure.

The coefficient limits have the untwisted transition law of a form
on $F$.  More explicitly, let $v=T(x)w$ be a holomorphic change of
bundle coordinates and let $t'=gt$.  The original matrices satisfy
\[
H'_m=|g|^{2/m}T^*H_mT.
\]
On a compact subset $K$ of the overlap, the error obtained by deleting
the scalar factor has, against a smooth compactly supported matrix
test function $\Phi$, the bound
\[
\left|\left\langle
\bigl(|g|^{2/m}-1\bigr)T^*H_mT,\Phi
\right\rangle\right|
\le C_K\|\Phi\|_{L^1}
\sup_K\bigl||g|^{2/m}-1\bigr|\longrightarrow0.
\]
Here the local matrix bound and the boundedness of $T$ are absorbed
into $C_K$.  Multiplication by the fixed smooth matrices $T,T^*$
preserves distributional convergence, as does a change of base
coordinates.  Consequently $H'=T^*HT$ almost everywhere.

\smallskip\noindent\textbf{Nonvanishing of the measurable limit.}
The limiting measurable form is not zero.  Choose maximizing unit
vectors $(x_m,v_m)$ in \eqref{eq:limit-normalization}.  After a further
subsequence they converge in the compact unit sphere bundle to
$(x_*,v_*)$.  This further subsequence retains all the distributional
limits already obtained.  In one fixed coordinate and frame chart
around $x_*$, let $c_m\to c$ be their fiber coordinates.  Then $c\ne0$,
and a local operator bound $\|H_m\|\le C$ gives
\[
\bigl|c^*H_m(x_m)c-c_m^*H_m(x_m)c_m\bigr|
\le C|c-c_m|(|c|+|c_m|)\longrightarrow0.
\]
On the other hand,
\[
c_m^*H_m(x_m)c_m=|t(x_m)|_0^{2/m}\longrightarrow1,
\]
since the norm of the nonvanishing frame is bounded above and away
from zero on a fixed compact neighborhood.  Thus the nonnegative
plurisubharmonic function
$u_m(x)=c^*H_m(x)c$ satisfies $u_m(x_m)\ge\tfrac12$ for large $m$.
Choose $\varrho>0$ so that $B(x_*,4\varrho)$ lies in the coordinate
chart, and choose a fixed smooth nonnegative function $\chi$ supported
in $B(x_*,3\varrho)$ and equal to one on $B(x_*,2\varrho)$.
For all sufficiently large $m$, the ball $B(x_m,\varrho)$ lies in
that latter ball.  The submean inequality yields
\begin{equation}\label{eq:limit-noncollapse}
\int\chi(x)c^*H_m(x)c\,dx
\ge\int_{B(x_m,\varrho)}u_m(x)\,dx
\ge\frac12\operatorname{Vol}(B(0,\varrho)).
\end{equation}
The left side is a pairing against the fixed smooth matrix test
function $\chi cc^*$, so the inequality passes to $H$.  It proves
that $H$ is nonzero as a measurable form.

\smallskip\noindent\textbf{A quadratic psh representative on every fiber.}
The measurable form is defined only up to a set of measure zero in
the base. Even a
pointwise limit of the traces may fail to be plurisubharmonic: in one
base variable, tracing $(1+|y|^2)|v|^2$ through $x=y^m$ gives
\[
 (1+|x|^{2/m})|v|^2.
\]
Its pointwise limit has coefficient $1$ at $x=0$ and $2$ elsewhere,
whereas its plurisubharmonic representative has coefficient $2$
everywhere. We now recover the intrinsic plurisubharmonic
representative of our measurable form and show that it remains a
Hermitian quadratic on every fiber.

Convolution regularization of
negative singular Hermitian metrics appears in
\cite[Section~3, Proposition~3.1]{BerndtssonPaun2008} and
\cite[Proposition~1.3(i)]{Raufi2015}; the use of polarization to recover
an everywhere-defined representative is also discussed in
\cite[Remark~2.2.3]{PaunTakayama2018}.
We supply the argument for the present possibly rank-deficient limit,
including the compatibility of the representatives on overlaps.
Work in a coordinate and bundle
chart $U$, and let $\rho$ be a nonnegative smooth radial function on
$\CC^2$, supported in the unit ball and of integral one.  Put
\[
\rho_\delta(y)=\delta^{-4}\rho(y/\delta),\qquad
H_\delta=H*\rho_\delta,\qquad
q_\delta(x,v)=v^*H_\delta(x)v,
\]
on the interior domain where this convolution is defined.
Convolving $H_m$ in the same way gives
\[
q_{m,\delta}(x,v)
=\int q_{m,t}(x-y,v)\rho_\delta(y)\,dy.
\]
It is plurisubharmonic jointly in $(x,v)$, since each base translation
$(x,v)\mapsto(x-y,v)$ is holomorphic and the average has nonnegative
weights.  For fixed $\delta$, distributional convergence of $H_m$
implies convergence of the convolved coefficients and all their
derivatives at every interior point: these are pairings against
translates of derivatives of the fixed smooth kernel.  The Levi
matrices of $q_{m,\delta}$ therefore converge pointwise to that of
$q_\delta$.  Hence $q_\delta$ is smooth and plurisubharmonic, and on
each smaller chart
\begin{equation}\label{eq:limit-convolution-bound}
0\le q_\delta(x,v)\le C|v|^2
\end{equation}
for all sufficiently small $\delta$.

These quadratics decrease as $\delta$ decreases.  To check this for
any fixed nonnegative radial kernel, write $\rho(y)=\phi(|y|)$ and
let $M_u(x,s)$ denote the spherical mean of a subharmonic function
$u$ on $\RR^4$.  Polar coordinates give
\begin{equation}\label{eq:limit-radial-means}
(u*\rho_\delta)(x)
=|\mathbb S^3|\int_0^1
M_u(x,\delta s)\phi(s)s^3\,ds.
\end{equation}
The spherical mean is nondecreasing in its radius.  The integration
weight in \eqref{eq:limit-radial-means} is nonnegative and does not
depend on $\delta$, so the convolution is nondecreasing in $\delta$.
Apply this to $u(x)=q_{m,t}(x,v)$ for each constant $v$, then pass to
the limit in $m$ at two fixed radii.  It follows that
\[
H_{\delta_1}(x)\le H_{\delta_2}(x)
\qquad(0<\delta_1<\delta_2)
\]
in the order of Hermitian forms.

Thus $q_\delta$ decreases, as $\delta\downarrow0$, to a finite
nonnegative plurisubharmonic function $q_U$ on the total-space chart.
The decreasing-limit assertion follows also by passing to the
submean inequalities; see \cite[Chapter~I, Theorem~5.4]{DemaillyCADG}.
Finiteness is supplied by \eqref{eq:limit-convolution-bound}.
At every fixed base point $x$, this limit is a Hermitian quadratic.
Indeed, all entries of $H_\delta(x)$ are determined by the finite
set of values at $e_j$, $e_j+e_k$, and $e_j+i e_k$, through
\begin{align*}
2\operatorname{Re}(H_\delta)_{jk}
&=q_\delta(x,e_j+e_k)-q_\delta(x,e_j)-q_\delta(x,e_k),\\
2\operatorname{Im}(H_\delta)_{jk}
&=q_\delta(x,e_j)+q_\delta(x,e_k)-q_\delta(x,e_j+i e_k).
\end{align*}
Each value on the right has a finite limit.  Hence the entries
converge to one Hermitian matrix $\widehat H_U(x)$, and passage to the
finite quadratic polynomial gives
\[
q_U(x,v)=v^*\widehat H_U(x)v\qquad\text{for every }v\in\CC^r.
\]
This matrix is positive semidefinite at every $x$, by nonnegativity.
At almost every $x$, approximate-identity convergence also gives
$\widehat H_U(x)=H(x)$.

Finally these local representatives agree on overlaps.  After the
holomorphic change of bundle coordinates, their measurable matrices
agree outside a base-null set, by the transition law already proved.
The quadratic functions therefore agree for every fiber vector
outside that set, hence almost everywhere in the total-space chart
by Fubini's theorem.  Both functions are locally bounded and
plurisubharmonic.  Such a function is determined everywhere by its
almost-everywhere values: the submean inequality and upper
semicontinuity give, in local Euclidean coordinates,
\[
u(p)=\lim_{\varepsilon\downarrow0}
\frac{1}{\operatorname{Vol}(B(p,\varepsilon))}
\int_{B(p,\varepsilon)}u.
\]
Thus the local $q_U$ glue to a single function $q$ on $F$.  Its
coefficients represent the original measurable form almost
everywhere, so \eqref{eq:limit-noncollapse} implies $q\not\equiv0$.
All the required pointwise and plurisubharmonic properties now hold.
\end{proof}

\begin{remark}\label{rmk:canonical-limit}
The function in Proposition~\ref{prop:limit} represents a
distributional limit; no convergence of restrictions to a divisor is
asserted. The graph identities will instead be evaluated directly in
this single everywhere-defined function $q$.
\end{remark}
